Uma \emph{função}~$f\colon A \to B$ mapeia todo elemento do
\emph{domínio}~$A$ em um único elemento do \emph{contradomínio}~$B$. Se
$x \in A$, chamamos o $y$ no qual~$f$ mapeia $x$ de
\emph{valor}~$f(x)$ de $f$ para o \emph{argumento}~$x$. Se $A$ é um conjunto de
pares, podemos pensar na função~$f$ como tomando dois argumentos. A
\emph{imagem}~$\ran{f}$ de~$f$ é o subconjunto de~$B$ que consiste em
todos os valores de~$f$.

Se $\ran{f} = B$, então $f$ é chamada de \emph{sobrejetora}. O valor~$f(x)$
é único no sentido de que $f$ mapeia $x$ em apenas um $f(x)$, nunca mais de um.
Se $f(x)$ também é único no sentido de que dois argumentos diferentes não
são mapeados no mesmo valor, $f$ é chamada de \emph{injetora}. As funções
que são tanto injetoras quanto sobrejetoras são chamadas de \emph{bijetoras}.

As funções bijetoras têm uma \emph{função inversa}~$f^{-1}$ única.
As funções também podem ser encadeadas: a função $(g \circ f)$ é
a \emph{composição} de $f$ com $g$. Composições de funções
injetoras são injetoras, e de funções sobrejetoras são sobrejetoras,
e $(f^{-1} \circ f)$ é a função identidade.

Se relaxamos a exigência de que $f$ deva ter um valor para todo $x
\in A$, obtemos a noção de \emph{função parcial}. Se $f\colon
A \pto B$ é parcial, dizemos que $f(x)$ é \emph{definida},
$f(x) \fdefined$, se $f$ tem um valor para o argumento~$x$, e, caso contrário,
dizemos que $f(x)$ é \emph{indefinida}, $f(x) \fundefined$. Toda função
(parcial)~$f$ está associada ao \emph{grafo}~$R_f$ de $f$, a
relação que vale sse $f(x) = y$.
